\documentclass[a5paper,10pt,dvipdfmx]{jsarticle}
\usepackage{amsmath,amssymb}
\usepackage[margin=17mm,headsep=5mm,includehead]{geometry}
\usepackage{fancyhdr,needspace,array,graphicx}
\pagestyle{fancy}
\fancyhf{}
\fancyhead[L]{\small 数学原論}
\fancyhead[R]{\small 記号・記法一覧\quad\thepage}
\setlength{\headheight}{14pt}
\renewcommand{\headrulewidth}{0pt}
\setlength{\parindent}{0pt}
\usepackage[dvipdfmx,hidelinks]{hyperref}
\usepackage{pxjahyper}
\newcommand{\NotationRef}[2]{\href{https://hawkoods.github.io/genron/?id=#1}{#2}}
\newsavebox{\NotationSymbol}
\newcommand{\NotationFit}[1]{\sbox{\NotationSymbol}{#1}%
  \ifdim\wd\NotationSymbol>\linewidth
    \resizebox{\linewidth}{!}{\usebox{\NotationSymbol}}%
  \else\usebox{\NotationSymbol}\fi}
\newcommand{\NotationRow}[2]{%
  \Needspace{3\baselineskip}\par\noindent
  \begin{tabular}{@{}>{\raggedright\arraybackslash}p{.43\linewidth}@{\hspace{.04\linewidth}}>{\raggedright\arraybackslash}p{.53\linewidth}@{}}
    \NotationFit{#1} & #2
  \end{tabular}\par\smallskip}
\newcommand{\NotationHeading}{%
  \Needspace{5\baselineskip}\par\noindent
  \begin{tabular}{@{}p{.43\linewidth}@{\hspace{.04\linewidth}}p{.53\linewidth}@{}}
    \textbf{記号} & \textbf{対象・意味}\\[2pt]\hline
  \end{tabular}\par\smallskip}
\begin{document}
\begin{center}
{\small 数学原論}\par\medskip
{\LARGE 記号・記法一覧}\par\smallskip
{\small 第2版}
\end{center}
記号とその指示対象を一行ずつ掲げる。参照付き項目は既刊ノートの知識台帳に対応する。参照のない項目は先行記法であり、定義の確立を意味しない。同一文脈で衝突する記号は添字又は説明で区別する。
\section{数学的言語・集合}
\NotationHeading
\NotationRow{\NotationRef{K-000002}{$\forall$}}{\NotationRef{K-000002}{全称量化}}
\NotationRow{\NotationRef{K-000002}{$\exists$}}{\NotationRef{K-000002}{存在量化}}
\NotationRow{\NotationRef{K-000006}{$\exists!$}}{\NotationRef{K-000006}{一意存在}}
\NotationRow{\NotationRef{K-000001}{$\neg$}}{\NotationRef{K-000001}{否定}}
\NotationRow{\NotationRef{K-000001}{$\land$}}{\NotationRef{K-000001}{連言}}
\NotationRow{\NotationRef{K-000001}{$\lor$}}{\NotationRef{K-000001}{選言}}
\NotationRow{\NotationRef{K-000001}{$\Rightarrow$}}{\NotationRef{K-000001}{含意}}
\NotationRow{\NotationRef{K-000001}{$\Leftrightarrow$}}{\NotationRef{K-000001}{同値}}
\NotationRow{\NotationRef{K-000021}{$\varnothing$}}{\NotationRef{K-000021}{空集合}}
\NotationRow{\NotationRef{K-000018}{$x\in A$}}{\NotationRef{K-000018}{所属}}
\NotationRow{\NotationRef{K-000018}{$x\notin A$}}{\NotationRef{K-000018}{非所属}}
\NotationRow{\NotationRef{K-000019}{$A\subseteq B$}}{\NotationRef{K-000019}{部分集合}}
\NotationRow{\NotationRef{K-000019}{$A\subsetneq B$}}{\NotationRef{K-000019}{真部分集合}}
\NotationRow{\NotationRef{K-000027}{$\mathcal P(X)$}}{\NotationRef{K-000027}{冪集合}}
\NotationRow{\NotationRef{K-000024}{$\{x\in X\mid P(x)\}$}}{\NotationRef{K-000024}{部分集合の指定}}
\NotationRow{\NotationRef{K-000029}{$A\cup B$}}{\NotationRef{K-000029}{合併}}
\NotationRow{\NotationRef{K-000029}{$A\cap B$}}{\NotationRef{K-000029}{共通部分}}
\NotationRow{\NotationRef{K-000029}{$A\setminus B$}}{\NotationRef{K-000029}{差集合}}
\NotationRow{\NotationRef{K-000044}{$(A_i)_{i\in I}$}}{\NotationRef{K-000044}{集合族}}
\NotationRow{\NotationRef{K-000045}{$\bigcup_{i\in I}A_i$}}{\NotationRef{K-000045}{集合族の合併}}
\NotationRow{\NotationRef{K-000045}{$\bigcap_{i\in I}A_i$}}{\NotationRef{K-000045}{集合族の共通部分}}
\NotationRow{\NotationRef{K-000033}{$A\times B$}}{\NotationRef{K-000033}{直積}}
\NotationRow{\NotationRef{K-000047}{$\prod_{i\in I}A_i$}}{\NotationRef{K-000047}{直積族}}
\NotationRow{\NotationRef{K-000049}{$\coprod_{i\in I}A_i$}}{\NotationRef{K-000049}{非交和}}
\NotationRow{\NotationRef{K-000032}{$(a,b)$}}{\NotationRef{K-000032}{順序対}}
\NotationRow{\NotationRef{K-000034}{$f:X\to Y$}}{\NotationRef{K-000034}{写像}}
\NotationRow{\NotationRef{K-000035}{$x\mapsto f(x)$}}{\NotationRef{K-000035}{写像の指定}}
\NotationRow{\NotationRef{K-000037}{$f(A)$}}{\NotationRef{K-000037}{像}}
\NotationRow{\NotationRef{K-000037}{$f^{-1}(B)$}}{\NotationRef{K-000037}{逆像}}
\NotationRow{\NotationRef{K-000041}{$f^{-1}$}}{\NotationRef{K-000041}{逆写像}}
\NotationRow{\NotationRef{K-000042}{$f\vert_A$}}{\NotationRef{K-000042}{制限}}
\NotationRow{\NotationRef{K-000036}{$g\circ f$}}{\NotationRef{K-000036}{合成}}
\NotationRow{\NotationRef{K-000036}{$\mathrm{id}_X$}}{\NotationRef{K-000036}{恒等写像}}
\NotationRow{\NotationRef{K-000047}{$Y^X$}}{\NotationRef{K-000047}{写像全体}}
\NotationRow{\NotationRef{K-000052}{$x\sim y$}}{\NotationRef{K-000052}{同値関係}}
\NotationRow{\NotationRef{K-000053}{$[x]$}}{\NotationRef{K-000053}{同値類}}
\NotationRow{\NotationRef{K-000053}{$X/{\sim}$}}{\NotationRef{K-000053}{商集合}}
\NotationRow{\NotationRef{K-000053}{$\pi$}}{\NotationRef{K-000053}{標準全射}}
\NotationRow{\NotationRef{K-000111}{$\lvert X\rvert$}}{\NotationRef{K-000111}{有限集合の元の個数}}
\NotationRow{\NotationRef{K-000099}{$\mathbb N$}}{\NotationRef{K-000099}{自然数集合（0を含む）}}
\NotationRow{\NotationRef{K-000068}{$\leq$}}{\NotationRef{K-000068}{半順序}}
\NotationRow{\NotationRef{K-000069}{$<$}}{\NotationRef{K-000069}{狭義順序}}
\NotationRow{\NotationRef{K-000073}{$\sup$}}{\NotationRef{K-000073}{上限}}
\NotationRow{\NotationRef{K-000073}{$\inf$}}{\NotationRef{K-000073}{下限}}
\NotationRow{\NotationRef{K-000079}{$\vee$}}{\NotationRef{K-000079}{結び}}
\NotationRow{\NotationRef{K-000079}{$\wedge$}}{\NotationRef{K-000079}{交わり}}
\NotationRow{\NotationRef{K-000070}{$P^{\mathrm{op}}$}}{\NotationRef{K-000070}{双対順序}}
\NotationRow{\NotationRef{K-000091}{$\operatorname{Fix}(c)$}}{\NotationRef{K-000091}{閉包作用素の固定元集合}}
\NotationRow{\NotationRef{K-000091}{$c$}}{\NotationRef{K-000091}{順序集合の閉包作用素}}
\NotationRow{\NotationRef{K-000093}{$f\dashv g$}}{\NotationRef{K-000093}{随伴（順序保存型）}}
\NotationRow{\NotationRef{K-000090}{$x^c$}}{\NotationRef{K-000090}{Boolean束の補元}}
\NotationRow{$\aleph_0$}{可算無限濃度\, {\scriptsize（先行）}}
\NotationRow{$\operatorname{card}X$}{一般の濃度\, {\scriptsize（先行）}}
\NotationRow{$\omega$}{最小無限順序数\, {\scriptsize（先行）}}
\section{位相・収束}
\NotationHeading
\NotationRow{\NotationRef{K-000146}{$\mathcal O$}}{\NotationRef{K-000146}{開集合系}}
\NotationRow{\NotationRef{K-000146}{$(X,\mathcal O)$}}{\NotationRef{K-000146}{位相空間}}
\NotationRow{\NotationRef{K-000146}{$U\in\mathcal O$}}{\NotationRef{K-000146}{開集合}}
\NotationRow{\NotationRef{K-000148}{$X\setminus U$}}{\NotationRef{K-000148}{開集合の補集合（閉集合）}}
\NotationRow{\NotationRef{K-000149}{$\mathcal C$}}{\NotationRef{K-000149}{閉集合系}}
\NotationRow{\NotationRef{K-000154}{$\mathcal O(\mathcal S)$}}{\NotationRef{K-000154}{生成位相}}
\NotationRow{\NotationRef{K-000146}{$\mathcal O_X$}}{\NotationRef{K-000146}{空間Xの位相}}
\NotationRow{\NotationRef{K-000152}{$\mathcal O_1\subseteq\mathcal O_2$}}{\NotationRef{K-000152}{位相の粗密}}
\NotationRow{$\mathcal N(x)$}{近傍系\, {\scriptsize（先行）}}
\NotationRow{$\mathcal B_x$}{近傍基\, {\scriptsize（先行）}}
\NotationRow{$\mathcal B$}{開基\, {\scriptsize（先行）}}
\NotationRow{$\mathcal S$}{準開基\, {\scriptsize（先行）}}
\NotationRow{$\operatorname{Int}(A)$}{内部\, {\scriptsize（先行）}}
\NotationRow{$\overline A$}{閉包\, {\scriptsize（先行）}}
\NotationRow{$\partial A$}{境界\, {\scriptsize（先行）}}
\NotationRow{$A'$}{集積点集合\, {\scriptsize（先行）}}
\NotationRow{$\operatorname{Ext}(A)$}{外部\, {\scriptsize（先行）}}
\NotationRow{$\mathcal O\vert_A$}{部分空間の位相\, {\scriptsize（先行）}}
\NotationRow{$d(x,y)$}{距離\, {\scriptsize（先行）}}
\NotationRow{$B_d(x,r)$}{開球\, {\scriptsize（先行）}}
\NotationRow{$\overline B_d(x,r)$}{閉球\, {\scriptsize（先行）}}
\NotationRow{$\operatorname{diam}A$}{直径\, {\scriptsize（先行）}}
\NotationRow{$d(x,A)$}{点と集合の距離\, {\scriptsize（先行）}}
\NotationRow{$\mathcal F$}{フィルター\, {\scriptsize（先行）}}
\NotationRow{$x_i\to x$}{収束\, {\scriptsize（先行）}}
\NotationRow{$\widehat X$}{完備化\, {\scriptsize（先行）}}
\NotationRow{$X\cong Y$}{同相\, {\scriptsize（先行）}}
\NotationRow{$X/{\sim}$}{商位相空間\, {\scriptsize（先行）}}
\NotationRow{$\pi_0(X)$}{経路連結成分の集合\, {\scriptsize（先行）}}
\NotationRow{$\pi_1(X,x)$}{基本群\, {\scriptsize（先行）}}
\NotationRow{$\pi_n(X,x)$}{高次ホモトピー群\, {\scriptsize（先行）}}
\NotationRow{$f\simeq g$}{ホモトピー\, {\scriptsize（先行）}}
\NotationRow{$\operatorname{supp}f$}{函数の台\, {\scriptsize（先行）}}
\section{代数・数体系}
\NotationHeading
\NotationRow{$\mathbb Z$}{整数環\, {\scriptsize（先行）}}
\NotationRow{$\mathbb Q$}{有理数体\, {\scriptsize（先行）}}
\NotationRow{$\mathbb R$}{実数体\, {\scriptsize（先行）}}
\NotationRow{$\mathbb C$}{複素数体\, {\scriptsize（先行）}}
\NotationRow{\NotationRef{K-000163}{$M$}}{\NotationRef{K-000163}{モノイド}}
\NotationRow{\NotationRef{K-000163}{$1_M$}}{\NotationRef{K-000163}{モノイドの単位元}}
\NotationRow{\NotationRef{K-000159}{$xy$}}{\NotationRef{K-000159}{二項演算の値}}
\NotationRow{\NotationRef{K-000168}{$x^n$}}{\NotationRef{K-000168}{自然数冪}}
\NotationRow{$G$}{群\, {\scriptsize（先行）}}
\NotationRow{$H\leq G$}{部分群\, {\scriptsize（先行）}}
\NotationRow{$N\trianglelefteq G$}{正規部分群\, {\scriptsize（先行）}}
\NotationRow{$1_G$}{群の単位元\, {\scriptsize（先行）}}
\NotationRow{$x^{-1}$}{群の逆元\, {\scriptsize（先行）}}
\NotationRow{$\langle S\rangle$}{生成部分群\, {\scriptsize（先行）}}
\NotationRow{$G/N$}{商群\, {\scriptsize（先行）}}
\NotationRow{$G/H$}{左剰余類集合\, {\scriptsize（先行）}}
\NotationRow{$[G:H]$}{部分群の指数\, {\scriptsize（先行）}}
\NotationRow{$Z(G)$}{群の中心\, {\scriptsize（先行）}}
\NotationRow{$C_G(S)$}{中心化群\, {\scriptsize（先行）}}
\NotationRow{$N_G(H)$}{正規化群\, {\scriptsize（先行）}}
\NotationRow{$[G,G]$}{導来部分群\, {\scriptsize（先行）}}
\NotationRow{$[x,y]$}{交換子\, {\scriptsize（先行）}}
\NotationRow{$\operatorname{ord}(x)$}{元の位数\, {\scriptsize（先行）}}
\NotationRow{$\operatorname{Aut}(G)$}{自己同型群\, {\scriptsize（先行）}}
\NotationRow{$\operatorname{Inn}(G)$}{内部自己同型群\, {\scriptsize（先行）}}
\NotationRow{$\ker f$}{準同型の核\, {\scriptsize（先行）}}
\NotationRow{$\operatorname{Im}f$}{準同型の像\, {\scriptsize（先行）}}
\NotationRow{$A,\ R$}{環（変数）\, {\scriptsize（先行）}}
\NotationRow{$K,\ k$}{体（変数）\, {\scriptsize（先行）}}
\NotationRow{$L/K$}{体の拡大\, {\scriptsize（先行）}}
\NotationRow{$[L:K]$}{拡大次数\, {\scriptsize（先行）}}
\NotationRow{$\mathfrak a$}{イデアル\, {\scriptsize（先行）}}
\NotationRow{$\mathfrak p$}{素イデアル\, {\scriptsize（先行）}}
\NotationRow{$\mathfrak m$}{極大イデアル\, {\scriptsize（先行）}}
\NotationRow{$(0)$}{零イデアル\, {\scriptsize（先行）}}
\NotationRow{$(a_1,\ldots,a_n)$}{生成イデアル\, {\scriptsize（先行）}}
\NotationRow{$A^\times$}{単元群\, {\scriptsize（先行）}}
\NotationRow{$A/\mathfrak a$}{商環\, {\scriptsize（先行）}}
\NotationRow{$\sqrt{\mathfrak a}$}{イデアルの根基\, {\scriptsize（先行）}}
\NotationRow{$\operatorname{Nil}(A)$}{冪零根基\, {\scriptsize（先行）}}
\NotationRow{$\operatorname{Jac}(A)$}{Jacobson根基\, {\scriptsize（先行）}}
\NotationRow{$\operatorname{Spec}A$}{素スペクトル\, {\scriptsize（先行）}}
\NotationRow{$\operatorname{MaxSpec}A$}{極大スペクトル\, {\scriptsize（先行）}}
\NotationRow{$S^{-1}A$}{局所化\, {\scriptsize（先行）}}
\NotationRow{$A_{\mathfrak p}$}{素イデアルでの局所化\, {\scriptsize（先行）}}
\NotationRow{$\operatorname{Frac}(A)$}{分数体\, {\scriptsize（先行）}}
\NotationRow{$A[X_1,\ldots,X_n]$}{多項式環\, {\scriptsize（先行）}}
\NotationRow{$A[[X]]$}{形式冪級数環\, {\scriptsize（先行）}}
\NotationRow{$\deg f$}{多項式の次数\, {\scriptsize（先行）}}
\NotationRow{$\operatorname{char}K$}{体の標数\, {\scriptsize（先行）}}
\NotationRow{$\operatorname{Gal}(L/K)$}{Galois群\, {\scriptsize（先行）}}
\NotationRow{$\overline K$}{代数閉包\, {\scriptsize（先行）}}
\NotationRow{$\mathcal O_K$}{代数体の整数環\, {\scriptsize（先行）}}
\NotationRow{$\operatorname{Tr}_{L/K}$}{体のトレース\, {\scriptsize（先行）}}
\NotationRow{$N_{L/K}$}{体のノルム\, {\scriptsize（先行）}}
\NotationRow{$v$}{付値\, {\scriptsize（先行）}}
\NotationRow{$\lvert\cdot\rvert_v$}{付値に対応する絶対値\, {\scriptsize（先行）}}
\NotationRow{$K_v$}{付値による完備化\, {\scriptsize（先行）}}
\NotationRow{$\mathbb Z_p$}{p進整数環\, {\scriptsize（先行）}}
\NotationRow{$\mathbb Q_p$}{p進数体\, {\scriptsize（先行）}}
\NotationRow{$\operatorname{Hom}_A(M,N)$}{加群準同型全体\, {\scriptsize（先行）}}
\NotationRow{$\operatorname{End}_A(M)$}{加群自己準同型環\, {\scriptsize（先行）}}
\NotationRow{$M\oplus N$}{直和\, {\scriptsize（先行）}}
\NotationRow{$M\otimes_A N$}{テンソル積\, {\scriptsize（先行）}}
\NotationRow{$M/\mathfrak a M$}{商加群\, {\scriptsize（先行）}}
\NotationRow{$M^*$}{代数的双対\, {\scriptsize（先行）}}
\section{線形代数・行列}
\NotationHeading
\NotationRow{$V,\ W$}{線形空間（変数）\, {\scriptsize（先行）}}
\NotationRow{$\dim_K V$}{次元\, {\scriptsize（先行）}}
\NotationRow{$\operatorname{span}_K S$}{線形包\, {\scriptsize（先行）}}
\NotationRow{$\mathcal L_K(V,W)$}{線形写像全体（純代数）\, {\scriptsize（先行）}}
\NotationRow{$V^*$}{双対空間\, {\scriptsize（先行）}}
\NotationRow{$V^{**}$}{二重双対空間\, {\scriptsize（先行）}}
\NotationRow{$T^\vee$}{双対写像\, {\scriptsize（先行）}}
\NotationRow{$A=(a_{ij})$}{行列\, {\scriptsize（先行）}}
\NotationRow{$M_{m,n}(K)$}{m行n列行列全体\, {\scriptsize（先行）}}
\NotationRow{$M_n(K)$}{n次正方行列全体\, {\scriptsize（先行）}}
\NotationRow{$I_n$}{単位行列\, {\scriptsize（先行）}}
\NotationRow{$0_{m,n}$}{零行列\, {\scriptsize（先行）}}
\NotationRow{$E_{ij}$}{行列単位\, {\scriptsize（先行）}}
\NotationRow{$A^{\mathsf T}$}{転置行列\, {\scriptsize（先行）}}
\NotationRow{$\overline A$}{成分ごとの複素共役行列\, {\scriptsize（先行）}}
\NotationRow{$A^*$}{共役転置行列\, {\scriptsize（先行）}}
\NotationRow{$A^{-1}$}{逆行列\, {\scriptsize（先行）}}
\NotationRow{$\det A$}{行列式\, {\scriptsize（先行）}}
\NotationRow{$\operatorname{tr}A$}{跡\, {\scriptsize（先行）}}
\NotationRow{$\operatorname{rank}A$}{階数\, {\scriptsize（先行）}}
\NotationRow{$\operatorname{adj}A$}{余因子転置行列\, {\scriptsize（先行）}}
\NotationRow{$\operatorname{diag}(a_1,\ldots,a_n)$}{対角行列\, {\scriptsize（先行）}}
\NotationRow{$[T]_{\mathcal B,\mathcal C}$}{線形写像の行列（始域基底B、終域基底C）\, {\scriptsize（先行）}}
\NotationRow{$[v]_{\mathcal B}$}{基底Bに関する座標列\, {\scriptsize（先行）}}
\NotationRow{$\chi_A(t)$}{特性多項式（det(tI-A)）\, {\scriptsize（先行）}}
\NotationRow{$m_A(t)$}{最小多項式\, {\scriptsize（先行）}}
\NotationRow{$\sigma(A)$}{スペクトル\, {\scriptsize（先行）}}
\NotationRow{$\lambda$}{固有値（変数）\, {\scriptsize（先行）}}
\NotationRow{$\ker(A-\lambda I)$}{固有空間\, {\scriptsize（先行）}}
\NotationRow{$\ker((A-\lambda I)^k)$}{一般化固有空間（十分大きいk）\, {\scriptsize（先行）}}
\NotationRow{$s_j(A)$}{特異値（降順）\, {\scriptsize（先行）}}
\NotationRow{$\mathrm{GL}_n(K)$}{一般線形群\, {\scriptsize（先行）}}
\NotationRow{$\mathrm{SL}_n(K)$}{特殊線形群\, {\scriptsize（先行）}}
\NotationRow{$\mathrm O(n)$}{直交群\, {\scriptsize（先行）}}
\NotationRow{$\mathrm U(n)$}{ユニタリ群\, {\scriptsize（先行）}}
\NotationRow{$\bigwedge^k V$}{外冪\, {\scriptsize（先行）}}
\NotationRow{$\operatorname{Sym}^k V$}{対称冪\, {\scriptsize（先行）}}
\NotationRow{$T(V)$}{テンソル代数\, {\scriptsize（先行）}}
\section{微分積分・解析}
\NotationHeading
\NotationRow{$f'(x)$}{一変数の導函数の値\, {\scriptsize（先行）}}
\NotationRow{$f^{(k)}(x)$}{k階導函数の値\, {\scriptsize（先行）}}
\NotationRow{$Df(x)$}{全微分（線形写像）\, {\scriptsize（先行）}}
\NotationRow{$D^kf(x)$}{k階全微分（多重線形写像）\, {\scriptsize（先行）}}
\NotationRow{$\partial_i f$}{第i偏導函数\, {\scriptsize（先行）}}
\NotationRow{$\partial^\alpha f$}{多重指数による偏導函数\, {\scriptsize（先行）}}
\NotationRow{$J_f(x)$}{Jacobi行列\, {\scriptsize（先行）}}
\NotationRow{$\det J_f(x)$}{Jacobi行列式（正方の場合）\, {\scriptsize（先行）}}
\NotationRow{$H_f(x)$}{Hesse行列（スカラー値函数）\, {\scriptsize（先行）}}
\NotationRow{$\nabla f$}{勾配\, {\scriptsize（先行）}}
\NotationRow{$\operatorname{div}X$}{発散\, {\scriptsize（先行）}}
\NotationRow{$\operatorname{rot}X$}{回転（3次元）\, {\scriptsize（先行）}}
\NotationRow{$\Delta f$}{Laplace作用素の適用\, {\scriptsize（先行）}}
\NotationRow{$\alpha!$}{多重指数階乗\, {\scriptsize（先行）}}
\NotationRow{$x^\alpha$}{多重指数冪\, {\scriptsize（先行）}}
\NotationRow{$o(h)$}{Landauの小さいo\, {\scriptsize（先行）}}
\NotationRow{$O(h)$}{Landauの大きいO\, {\scriptsize（先行）}}
\NotationRow{$\limsup$}{上極限\, {\scriptsize（先行）}}
\NotationRow{$\liminf$}{下極限\, {\scriptsize（先行）}}
\NotationRow{$\sum_{n=0}^\infty a_n$}{級数\, {\scriptsize（先行）}}
\NotationRow{$\int_a^b f(x)\,dx$}{区間上の積分\, {\scriptsize（先行）}}
\NotationRow{$\operatorname{Re}z$}{実部\, {\scriptsize（先行）}}
\NotationRow{$\operatorname{Im}z$}{虚部\, {\scriptsize（先行）}}
\NotationRow{$\overline z$}{複素共役\, {\scriptsize（先行）}}
\NotationRow{$\lvert z\rvert$}{絶対値\, {\scriptsize（先行）}}
\NotationRow{$\arg z$}{偏角（枝を明示）\, {\scriptsize（先行）}}
\NotationRow{$\operatorname{Res}_{z=a}f$}{留数\, {\scriptsize（先行）}}
\NotationRow{$\oint_\gamma f(z)\,dz$}{閉曲線上の複素積分\, {\scriptsize（先行）}}
\NotationRow{$\mathcal H(U)$}{正則函数全体\, {\scriptsize（先行）}}
\NotationRow{$\widehat f$}{Fourier変換\, {\scriptsize（先行）}}
\NotationRow{$\mathcal F f$}{Fourier変換作用素の値\, {\scriptsize（先行）}}
\NotationRow{$f*g$}{畳込み\, {\scriptsize（先行）}}
\NotationRow{$\mathcal D(U)$}{試験函数空間\, {\scriptsize（先行）}}
\NotationRow{$\mathcal D'(U)$}{超函数空間\, {\scriptsize（先行）}}
\NotationRow{$\mathcal S(\mathbb R^n)$}{Schwartz空間\, {\scriptsize（先行）}}
\NotationRow{$\mathcal S'(\mathbb R^n)$}{緩増加超函数空間\, {\scriptsize（先行）}}
\NotationRow{$W^{k,p}(U)$}{Sobolev空間\, {\scriptsize（先行）}}
\NotationRow{$H^k(U)$}{L2型Sobolev空間\, {\scriptsize（先行）}}
\NotationRow{$\lVert x\rVert$}{ノルム\, {\scriptsize（先行）}}
\NotationRow{$\langle x,y\rangle$}{内積\, {\scriptsize（先行）}}
\NotationRow{$X^*$}{連続双対空間\, {\scriptsize（先行）}}
\NotationRow{$\mathcal L(X,Y)$}{有界線形作用素全体\, {\scriptsize（先行）}}
\NotationRow{$T^*$}{Hilbert空間の随伴作用素\, {\scriptsize（先行）}}
\NotationRow{$\lVert T\rVert_{\mathrm{op}}$}{作用素ノルム\, {\scriptsize（先行）}}
\NotationRow{$\sigma(T)$}{作用素のスペクトル\, {\scriptsize（先行）}}
\NotationRow{$\rho(T)$}{レゾルベント集合\, {\scriptsize（先行）}}
\NotationRow{$R(\lambda,T)$}{レゾルベント\, {\scriptsize（先行）}}
\NotationRow{$C(X)$}{連続函数全体\, {\scriptsize（先行）}}
\NotationRow{$C_c(X)$}{コンパクト台連続函数全体\, {\scriptsize（先行）}}
\NotationRow{$C_0(X)$}{無限遠で消える連続函数全体\, {\scriptsize（先行）}}
\NotationRow{$C^k(U)$}{k回連続微分可能函数全体\, {\scriptsize（先行）}}
\NotationRow{$C^\infty(U)$}{滑らかな函数全体\, {\scriptsize（先行）}}
\section{測度・確率}
\NotationHeading
\NotationRow{\NotationRef{K-000182}{$\mathcal A$}}{\NotationRef{K-000182}{σ-加法族}}
\NotationRow{$(X,\mathcal A)$}{可測空間\, {\scriptsize（先行）}}
\NotationRow{\NotationRef{K-000187}{$\sigma(\mathcal S)$}}{\NotationRef{K-000187}{生成σ-加法族}}
\NotationRow{$\mathcal B(X)$}{Borel加法族\, {\scriptsize（先行）}}
\NotationRow{$\mu$}{測度\, {\scriptsize（先行）}}
\NotationRow{$(X,\mathcal A,\mu)$}{測度空間\, {\scriptsize（先行）}}
\NotationRow{$\mu^*$}{外測度\, {\scriptsize（先行）}}
\NotationRow{$\mu\otimes\nu$}{積測度\, {\scriptsize（先行）}}
\NotationRow{$\mathbf 1_A$}{指示函数\, {\scriptsize（先行）}}
\NotationRow{$\int_A f\,d\mu$}{測度に関する積分\, {\scriptsize（先行）}}
\NotationRow{$\frac{d\nu}{d\mu}$}{Radon--Nikodym微分\, {\scriptsize（先行）}}
\NotationRow{$f_*\mu$}{像測度\, {\scriptsize（先行）}}
\NotationRow{$\mu\ll\nu$}{絶対連続\, {\scriptsize（先行）}}
\NotationRow{$\mu\perp\nu$}{相互特異\, {\scriptsize（先行）}}
\NotationRow{$L^p(X,\mu)$}{Lebesgue空間\, {\scriptsize（先行）}}
\NotationRow{$\lVert f\rVert_p$}{Lpノルム\, {\scriptsize（先行）}}
\NotationRow{$\ell^p$}{数列空間\, {\scriptsize（先行）}}
\NotationRow{$(\Omega,\mathcal F,\mathbb P)$}{確率空間\, {\scriptsize（先行）}}
\NotationRow{$\mathbb E[X]$}{期待値\, {\scriptsize（先行）}}
\NotationRow{$\mathbb E[X\mid\mathcal G]$}{条件付き期待値\, {\scriptsize（先行）}}
\NotationRow{$\operatorname{Var}(X)$}{分散\, {\scriptsize（先行）}}
\NotationRow{$\operatorname{Cov}(X,Y)$}{共分散\, {\scriptsize（先行）}}
\NotationRow{$\mathcal L(X)$}{確率変数の法則\, {\scriptsize（先行）}}
\NotationRow{$(\mathcal F_t)_{t\geq0}$}{フィルトレーション\, {\scriptsize（先行）}}
\NotationRow{$(X_t)_{t\geq0}$}{確率過程\, {\scriptsize（先行）}}
\NotationRow{$B_t$}{Brown運動\, {\scriptsize（先行）}}
\NotationRow{$X_n\xrightarrow{\mathrm{a.s.}}X$}{概収束\, {\scriptsize（先行）}}
\NotationRow{$X_n\xrightarrow{\mathbb P}X$}{確率収束\, {\scriptsize（先行）}}
\NotationRow{$X_n\xrightarrow{\mathrm d}X$}{分布収束\, {\scriptsize（先行）}}
\section{圏・幾何・ホモロジー}
\NotationHeading
\NotationRow{$\mathcal C,\mathcal D$}{圏（変数）\, {\scriptsize（先行）}}
\NotationRow{$\operatorname{Ob}(\mathcal C)$}{対象全体\, {\scriptsize（先行）}}
\NotationRow{$\operatorname{Hom}_{\mathcal C}(X,Y)$}{射集合\, {\scriptsize（先行）}}
\NotationRow{$F:\mathcal C\to\mathcal D$}{函手\, {\scriptsize（先行）}}
\NotationRow{$\mathcal C^{\mathrm{op}}$}{反対圏\, {\scriptsize（先行）}}
\NotationRow{$\eta:F\Rightarrow G$}{自然変換\, {\scriptsize（先行）}}
\NotationRow{$\varprojlim$}{射影極限\, {\scriptsize（先行）}}
\NotationRow{$\varinjlim$}{帰納極限\, {\scriptsize（先行）}}
\NotationRow{$F\dashv G$}{随伴函手\, {\scriptsize（先行）}}
\NotationRow{$(M,\mathcal A)$}{微分多様体\, {\scriptsize（先行）}}
\NotationRow{$(U,\varphi)$}{座標近傍\, {\scriptsize（先行）}}
\NotationRow{$\psi\circ\varphi^{-1}$}{座標変換\, {\scriptsize（先行）}}
\NotationRow{$T_pM$}{接空間\, {\scriptsize（先行）}}
\NotationRow{$TM$}{接束\, {\scriptsize（先行）}}
\NotationRow{$T_p^*M$}{余接空間\, {\scriptsize（先行）}}
\NotationRow{$T^*M$}{余接束\, {\scriptsize（先行）}}
\NotationRow{$df_p$}{接写像\, {\scriptsize（先行）}}
\NotationRow{$f^*\omega$}{微分形式の引戻し\, {\scriptsize（先行）}}
\NotationRow{$\mathfrak X(M)$}{ベクトル場全体\, {\scriptsize（先行）}}
\NotationRow{$\Omega^k(M)$}{微分k形式全体\, {\scriptsize（先行）}}
\NotationRow{$d\omega$}{外微分\, {\scriptsize（先行）}}
\NotationRow{$\omega\wedge\eta$}{外積\, {\scriptsize（先行）}}
\NotationRow{$\iota_X\omega$}{内部積\, {\scriptsize（先行）}}
\NotationRow{$\mathcal L_X\omega$}{Lie微分\, {\scriptsize（先行）}}
\NotationRow{$g$}{Riemann計量\, {\scriptsize（先行）}}
\NotationRow{$\nabla$}{接続\, {\scriptsize（先行）}}
\NotationRow{$\nabla_XY$}{共変微分\, {\scriptsize（先行）}}
\NotationRow{$R$}{曲率テンソル\, {\scriptsize（先行）}}
\NotationRow{$\operatorname{Ric}$}{Ricci曲率\, {\scriptsize（先行）}}
\NotationRow{$\operatorname{grad}f$}{Riemann勾配\, {\scriptsize（先行）}}
\NotationRow{$\operatorname{Hess}f$}{Riemann Hessian\, {\scriptsize（先行）}}
\NotationRow{$\mathfrak g$}{Lie代数\, {\scriptsize（先行）}}
\NotationRow{$[X,Y]$}{Lie括弧\, {\scriptsize（先行）}}
\NotationRow{$\exp$}{指数写像\, {\scriptsize（先行）}}
\NotationRow{$C_n(X)$}{鎖群\, {\scriptsize（先行）}}
\NotationRow{$\partial_n$}{境界写像\, {\scriptsize（先行）}}
\NotationRow{$Z_n(X)$}{サイクル群\, {\scriptsize（先行）}}
\NotationRow{$B_n(X)$}{境界群\, {\scriptsize（先行）}}
\NotationRow{$H_n(X;A)$}{ホモロジー群\, {\scriptsize（先行）}}
\NotationRow{$H^n(X;A)$}{コホモロジー群\, {\scriptsize（先行）}}
\NotationRow{$\operatorname{Ext}_A^n(M,N)$}{Ext函手\, {\scriptsize（先行）}}
\NotationRow{$\operatorname{Tor}_n^A(M,N)$}{Tor函手\, {\scriptsize（先行）}}
\NotationRow{$(X,\mathcal O_X)$}{環付き空間\, {\scriptsize（先行）}}
\NotationRow{$\mathcal F(U)$}{層の切断集合\, {\scriptsize（先行）}}
\NotationRow{$\mathcal F_x$}{層の茎\, {\scriptsize（先行）}}
\NotationRow{$V(I)$}{イデアルの零点集合\, {\scriptsize（先行）}}
\NotationRow{$D(f)$}{基本開集合\, {\scriptsize（先行）}}
\NotationRow{$\operatorname{Proj}S$}{射影スペクトル\, {\scriptsize（先行）}}
\NotationRow{$G=(V,E)$}{グラフ\, {\scriptsize（先行）}}
\NotationRow{$\deg(v)$}{頂点の次数\, {\scriptsize（先行）}}
\NotationRow{$A_G$}{隣接行列\, {\scriptsize（先行）}}
\NotationRow{$\operatorname*{argmin}_{x\in X}f(x)$}{最小化点集合\, {\scriptsize（先行）}}
\NotationRow{$\operatorname*{argmax}_{x\in X}f(x)$}{最大化点集合\, {\scriptsize（先行）}}
\section{調査記録}
調査記録：Bredon, Topology and Geometry, 近傍の定義2.3--2.5；Atiyah--Macdonald, Introduction to Commutative Algebra (1969), Notation and Terminology, Ch.1；Neukirch, Algebraic Number Theory (1999), Ch. I 第3；Mac Lane, Categories for the Working Mathematician, Ch. I 第1；Lee, Introduction to Smooth Manifolds, 2nd ed. (2013), Ch.1 の座標変換；Conway, A Course in Functional Analysis, Ch. I 第3・Ch. II 第2。Google Driveの抽出本文を確認。その他の先行記法は標準的慣習からの独自提案であり、文献の逐語転記・全分野の文献照合ではない。

第2版の追加記法は編者の先行提案である。転置は右上の立体T、共役転置は星印とする。Jacobi行列は全微分の標準基底に関する行列で、Hesse行列は二階全微分の行列表示である。複素内積の線形側、交換子の順序、Fourier変換の正規化、曲率・レゾルベントの符号は正式導入時に規定する。
\end{document}
